% ====================================================================
% §2 통계역학 기초 (Part 0) — 파트 a: 앙상블·단일자리·격자기체 (2.1–2.3)
% ====================================================================
\section{통계역학 기초 --- 분배함수에서 Nernst 식까지 (Part 0)}\label{sec:sm-found}

본론($\mathrm N1$--$\mathrm N9$)의 모든 평형식은 통계역학의 한 사슬에서 나온다. 이 절은 그 사슬을 처음부터 놓으며, 통계역학
사전지식 없이도 본론의 출발점들이 이 절 안에서 닫히도록 각 단계를 (a) 출발 $\to$ (b) 연산 $\to$ (c) 중간식
$\to$ (d) 박스로 전개한다. 유도의 골격은 고전 통계역학(분배함수$\cdot$화학퍼텐셜)으로 닫히고, 양자 통계가
필요한 항(점유 자리의 진동 자유도, Part T 의 포논$\cdot$전자 분포, \S\ref{sec:vibel})은 해당 자리에서 명시적으로 도입한다.
사슬은
\[
\boxed{\;\begin{aligned}
&\text{미시상태·앙상블}\ \to\ Z,\Xi\ \to\ \theta=\frac{1}{1+e^{\beta(\tilde\varepsilon-\mu)}}\ \to\ \text{lattice gas}\\
&\to\ \mu(\theta;\Omega)\ \to\ \xi_\eq\ \text{logistic}\ \to\ \textstyle\sum_jQ_j\,\xi_{\eq,j}=Q\bar x\ \to\ G,\ \text{Nernst}
\end{aligned}\;}
\]
이고, 도착점들이 본론 결과 사슬 --- \S\ref{sec:center} 의 $U_j(T)$, \S\ref{sec:hys} 의 히스테리시스,
\S\ref{sec:width} 의 $\xi_\eq$, \S\ref{sec:eqpeak} 의 전하 보존식 --- 의 입력이다.

\subsection{앙상블과 Boltzmann·Gibbs 인자 --- 분배함수와 화학퍼텐셜 $\mu$}\label{sec:sm-ensemble}
\textbf{(a) 출발 --- 미시상태·앙상블과 두 공준.} \emph{미시상태}(microstate)는 계의 모든 자유도를 지정한
배치 하나이고, \emph{거시상태}(macrostate)는 에너지·부피·입자수 $(E,V,N)$ 같은 소수의 거시변수만 지정한
미시상태의 집합이다. 같은 거시 조건을 공유하는 미시상태들의 모음이 \emph{앙상블}(ensemble)이다. 통계역학은
두 공준으로 시작한다 --- 고립계 평형에서 접근 가능한 $W$ 개 미시상태는 등확률이고(공준 1), 엔트로피는 그
경우의 수의 로그다(공준 2, Boltzmann):
\begin{equation}
S\;=\;k_B\ln W.
\label{eq:sm-S}
\end{equation}
열역학과의 접점은 근본 미분식 $\dd E=T\dd S-P\dd V+\mu\,\dd N$ 을 $S$ 에 대해 푼 형태다:
\begin{equation}
\dd S=\frac{1}{T}\,\dd E+\frac{P}{T}\,\dd V-\frac{\mu}{T}\,\dd N
\quad\Longrightarrow\quad
\frac{\partial S}{\partial E}\Big|_{V,N}=\frac1T,\qquad
\frac{\partial S}{\partial N}\Big|_{E,V}=-\frac{\mu}{T}.
\label{eq:sm-fund}
\end{equation}
여기서 $\mu$ 의 물리적 의미가 이미 보인다 --- 입자 1개를 받아들일 때 엔트로피가 $\mu/T$ 만큼 \emph{감소}하는
값비쌈의 척도다. 두 계 1,2 가 입자를 교환하면 총 엔트로피 변화가
$\dd S_\mathrm{tot}=\big(\tfrac{\mu_2}{T}-\tfrac{\mu_1}{T}\big)\dd N_1\ge0$ 이므로 입자는 $\mu$ 가 높은 쪽에서
낮은 쪽으로 흐르고, 정지 조건은 $\mu_1=\mu_2$ --- 화학퍼텐셜의 일치가 곧 입자 교환 평형이다.

\textbf{(b) 연산 --- 저장조와 접촉한 작은 계.} 관심 계가 거대한 저장조(reservoir)와 에너지·입자를 교환한다.
계가 미시상태 $i$(에너지 $E_i$, 입자수 $N_i$)에 있을 확률은 공준 1 에 의해 그때 저장조가 가질 수 있는 경우의
수에 비례한다:
\begin{equation}
p_i\;\propto\;W_R(E_\mathrm{tot}-E_i,\,N_\mathrm{tot}-N_i)
\;=\;\exp\!\Big[\frac{S_R(E_\mathrm{tot}-E_i,\,N_\mathrm{tot}-N_i)}{k_B}\Big].
\label{eq:sm-resv}
\end{equation}
저장조가 계보다 훨씬 크므로 $S_R$ 을 1차에서 절단해 전개하면(식~\eqref{eq:sm-fund} 의 두 도함수 대입)
\begin{equation}
S_R(E_\mathrm{tot}-E_i,\,N_\mathrm{tot}-N_i)
= S_R(E_\mathrm{tot},N_\mathrm{tot})-\frac{E_i}{T}+\frac{\mu N_i}{T}+\mathcal O(E_i^2,N_i^2).
\label{eq:sm-taylor}
\end{equation}
\textbf{(c) 중간식 --- Gibbs 인자와 canonical 특수형.} 식~\eqref{eq:sm-taylor} 을 \eqref{eq:sm-resv} 에 넣으면
상수 몫이 규격화로 빠지고
\begin{equation}
p_i\;\propto\;e^{-\beta(E_i-\mu N_i)},\qquad \beta\equiv\frac{1}{k_BT}
\label{eq:sm-gibbsfactor}
\end{equation}
--- 이 지수가 \emph{Gibbs 인자}다. 입자 교환이 없으면($N_i$ 고정) $e^{-\beta E_i}$ 의 \emph{Boltzmann 인자}로
줄고, 그 규격화 합이 canonical 분배함수(partition function) $Z=\sum_ie^{-\beta E_i}$, 그 로그가 Helmholtz
자유에너지 $F=-k_BT\ln Z$ 다($\langle E\rangle=\partial(\beta F)/\partial\beta$·$S=-\partial F/\partial T$ 가
열역학 관계와 일치함이 이 명명의 근거다). 이 $F$ 는 Faraday 상수 $F$ 와 기호만 같고 문맥으로 갈린다 ---
전자는 $\partial F/\partial T$ 처럼 \emph{미분당하는 자유에너지}, 후자는 $FV$·$RT/F$ 의 전하 환산
곱$\cdot$나눗셈 상수다. 식~\eqref{eq:sm-fund} 의 $\mu$ 를 $F$ 로 옮기면
$\mu=\partial F/\partial N|_{T,V}$ --- \emph{온도·부피를 고정한 채 입자 1개를 더 넣는 자유에너지 비용}이며,
이것이 본론 \S\ref{sec:center} 가 쓰는 거시 정의 $\mu=\partial G/\partial n|_{T,P}$ 와 같은 양임은
\S\ref{sec:sm-macro} 에서 닫는다.

\textbf{(d) 박스 --- grand canonical 분배함수.} 입자수까지 교환하는 앙상블이
\emph{대정준}(grand canonical) 앙상블이고, 그 규격화 합과 평균 입자수는
\begin{equation}
\boxed{\;\Xi=\sum_i e^{-\beta(E_i-\mu N_i)},\qquad
p_i=\frac{e^{-\beta(E_i-\mu N_i)}}{\Xi},\qquad
\langle N\rangle=k_BT\,\frac{\partial\ln\Xi}{\partial\mu}\;}
\label{eq:sm-gc}
\end{equation}
($\partial\ln\Xi/\partial\mu=\beta\sum_iN_ip_i=\beta\langle N\rangle$ 로 셋째 식이 첫째·둘째에서 나온다).
\begin{verifybox}
극한 검산 --- $\beta\to0$ 이면 $p_i\to$ 등확률(열적 규율 소실), $\beta\to\infty$ 면 $E_i-\mu N_i$ 최소 상태만
남는다. 삽입 전극이 정확히 이 설정이다 --- 격자의 자리들이 전해질을 통해 Li 저장조(상대극 Li 금속)와 입자를
교환하므로 $\mu$ 가 고정된 대정준 문제이고(그림~\ref{fig:sm-reservoir}), 다음 소절이 그 최소 단위(자리
하나)를 푼다.
\end{verifybox}

\begin{figure}[t]
\centering
\begin{tikzpicture}[font=\scriptsize,
  box/.style={draw,rounded corners=2pt,align=center,inner sep=4pt},
  ar/.style={{Latex[length=1.5mm]}-{Latex[length=1.5mm]},thick}]
\node[box,minimum width=30mm,minimum height=16mm] (sys) {system\\insertion sites\\microstate $i$: $(E_i,\,N_i)$};
\node[box,minimum width=46mm,minimum height=16mm,fill=black!6,right=20mm of sys] (res)
  {reservoir: $T,\ \mu$ fixed\\Li metal counter electrode\\$+$ electrolyte};
\draw[ar] ([yshift=3.2mm]sys.east) -- ([yshift=3.2mm]res.west) node[midway,above] {energy $E$};
\draw[ar] ([yshift=-3.2mm]sys.east) -- ([yshift=-3.2mm]res.west) node[midway,below] {particles $N$ (Li)};
\end{tikzpicture}
\caption{grand canonical 구도. 관심 계(삽입 자리들)는 상대극$\cdot$전해질이 이루는 큰 저장조와
에너지$\cdot$입자를 모두 교환하고, 저장조는 $T$ 와 $\mu$ 를 고정한다. 에너지만 열면 Boltzmann 인자, 입자까지
열면 Gibbs 인자(식~\eqref{eq:sm-gibbsfactor}) --- 지수의 자리가 하나 늘어날 뿐 논리는 같다.}
\label{fig:sm-reservoir}
\end{figure}

\subsection{단일 삽입 자리 --- 점유 통계와 fluctuation}\label{sec:sm-site}
앞 소절이 놓은 대정준 도구(식~\eqref{eq:sm-gc})를 가장 작은 계 --- 격자의 삽입 자리 \emph{하나} --- 에
적용한다. 자리는 전해질 너머의 Li 저장조와 입자를 교환하므로 식~\eqref{eq:sm-gc} 의 대정준 설정 그대로다.
전개는 두 단으로 놓는다 --- 먼저 이온을 내부 구조 없는 점입자로 이상화한 두-상태 원형을 표준 유도
그대로 닫고, 이어 점유 이온의 진동 자유도를 켠 확장이 같은 함수형으로 되돌아옴을 보인 뒤, 점유의
fluctuation 까지 읽는다. 도착점인 점유율 $\theta$ 와 유효 자리 자유에너지 $\tilde\varepsilon(T)$ 는
\S\ref{sec:sm-lattice} 의 lattice gas 와 \S\ref{sec:sm-electro} 의 전기화학 결선이 받는 입력이다.

\textbf{(a) 출발 --- 두 가정과 원형의 두-상태 목록.} 유도에 앞서 두 가정을 명시한다.
\begin{itemize}[leftmargin=2.2em,itemsep=0.1em]
\item \textbf{가정 1 (hard-core 배타).} 한 자리에는 Li 이온이 $0$ 개 또는 $1$ 개만 들어간다. 근거는 양자
통계가 아니라 강한 단거리 반발이다 --- 자리의 부피와 이온 사이 정전$\cdot$입체 반발이 이중 점유의
에너지를 사실상 무한대로 올린다. 이 가정이 점유수 합을 $n\in\{0,1\}$ 로 자른다.
\item \textbf{가정 2 (자리 독립).} 이 소절에서는 자리의 에너지가 이웃 자리의 점유와 무관하다고 둔다.
이 가정은 \S\ref{sec:sm-mf} 에서 상호작용 파라미터 $\Omega$ 로 완화된다.
\end{itemize}
원형에서는 여기에 이상화 하나를 더 얹는다 --- 이온을 내부 구조 없는 점입자로 두어, 점유 이온이 우물
안에서 갖는 위치$\cdot$운동량 배치를 잠시 얼린다(이 이상화는 아래 확장이 걷어낸다). 그러면 미시상태의
목록은 단 둘이다: 빈 자리($n=0$, 에너지 $0$ 기준)와 점유 자리($n=1$, 에너지 $\varepsilon_0$).

\textbf{(b) 연산 --- 두 항의 대정준 합.} 각 미시상태의 무게는 Gibbs 인자
$e^{-\beta(E_i-\mu N_i)}$(식~\eqref{eq:sm-gibbsfactor})이고, 목록이 둘뿐이라 대정준 합도 두 항이다 ---
빈 자리가 $1$, 점유 자리가 $e^{-\beta(\varepsilon_0-\mu)}$:
\begin{equation}
\Xi_1^{0}\;=\;\sum_{n=0,1}e^{-\beta(\varepsilon_0-\mu)\,n}
\;=\;1+e^{-\beta(\varepsilon_0-\mu)}.
\label{eq:sm-baresum}
\end{equation}
기호에 관하여 --- $\Xi$ 는 자리 1개의 \emph{대정준} 분배함수로, \S\ref{sec:sm-ensemble} 의 canonical
$Z$ 와는 앙상블이 다르므로 기호부터 구분해 적고(첨자 $1$ = 자리 1개), 위 첨자 $0$ 은 내부 자유도를 얼린
원형임을 표시한다(평균 점유도 같은 표기 $\langle n\rangle^{0}$ 로 구분한다).

\textbf{(c) 중간식 --- 상태 확률과 평균 점유.} 식~\eqref{eq:sm-gc} 의 둘째 식으로 두 상태의 확률은
$p_0=1/\Xi_1^{0}$·$p_1=e^{-\beta(\varepsilon_0-\mu)}/\Xi_1^{0}$ 이고, 평균 점유는 그 확률 가중 평균이다:
\begin{equation}
\langle n\rangle^{0}
\;=\;0\cdot p_0+1\cdot p_1
\;=\;\frac{e^{-\beta(\varepsilon_0-\mu)}}{1+e^{-\beta(\varepsilon_0-\mu)}},
\label{eq:sm-baremid}
\end{equation}
그리고 식~\eqref{eq:sm-gc} 의 셋째 식 $k_BT\,\partial\ln\Xi_1^{0}/\partial\mu$ 를
식~\eqref{eq:sm-baresum} 에 적용해도 같은 결과가 나온다(교차검증).

\textbf{(d) 박스 --- 두-상태 원형 $\Xi_1^{0}$·$\langle n\rangle^{0}$.} 식~\eqref{eq:sm-baremid} 의
분자$\cdot$분모를 $e^{-\beta(\varepsilon_0-\mu)}$ 로 나누면
\begin{equation}
\boxed{\;\Xi_1^{0}\;=\;1+e^{-\beta(\varepsilon_0-\mu)},\qquad
\langle n\rangle^{0}\;=\;\frac{1}{1+e^{+\beta(\varepsilon_0-\mu)}}\;.}
\label{eq:sm-bare}
\end{equation}
이 두-상태 대정준 유도는 국소화 흡착(localized adsorption)의 Langmuir 등온선 표준 유도
그대로이고 \cite{hill1960,fowler1939}, 삽입 전극에의 적용 원형은 McKinnon--Haering
\cite{mckinnon1983} 이다. 식~\eqref{eq:sm-bare} 의 $\langle n\rangle^{0}$ 은 준위(state) 하나의 평균 전자 점유를
주는 Fermi--Dirac 함수 $f(E)=1/(e^{\beta(E-\mu)}+1)$ 에서 $E$ 자리에 $\varepsilon_0$ 이 앉은 것과
\emph{동형}이다 --- 이 대응의 물리적 배경(페르미온$\cdot$보손의 양자 통계)은 별도의 배경 박스가
자족적으로 정리한다.
\begin{srcbox}
\textbf{다리 --- eq:sm-bare 와 eq:Veq 가 어느 삽입 등온선인가.} 식~\eqref{eq:sm-bare} 의 두-상태 원형과
\S\ref{sec:hys} 의 상호작용형 식~\eqref{eq:Veq} 는 McKinnon--Haering 이 삽입 전극에 세운 격자기체
등온선\cite{mckinnon1983} 의 원형$\cdot$확장에 각각 대응한다:
\begin{center}\footnotesize
\begin{tabular}{@{}ll@{}}
\hline
본문 & McKinnon--Haering\cite{mckinnon1983} 대응(격자기체) \\
\hline
평균 점유 $\langle n\rangle^{0}$ (식~\eqref{eq:sm-bare}) & 삽입 분율 $x$ (자리 점유율 --- 본문 $\theta$$\cdot$$\xi$ 에 해당) \\
$\tfrac{RT}{sF}\ln\tfrac{\xi}{1-\xi}$ 항 (식~\eqref{eq:Veq}) & 배치 항 $\tfrac{kT}{e}\ln\tfrac{x}{1-x}$ \\
표준 중심 $U_j$ & 표준 전위 $E_0$ \\
상호작용 몫 $\tfrac{\Omega_j}{sF}(1-2\xi)$ (식~\eqref{eq:Veq}) & 평균장 자리-자리 상호작용 항 \\
문턱 $\Omega_j>2RT$ (식~\eqref{eq:spinodal}) & 상호작용이 1차 전이(전위 계단)를 낳는 임계 세기 \\
\hline
\end{tabular}
\end{center}
\emph{방법 요지} --- McKinnon--Haering 은 삽입 자리를 동등한 격자기체로 보고 각 자리 점유를 대정준(Langmuir
형) 평형으로 닫아 삽입 등온선(전위--조성 관계)을 얻고, 자리 간 평균장 상호작용을 더해 그 세기가 임계를
넘으면 1차 전이(전위 계단$\cdot$plateau)가 나타남을 보인다. \emph{가정 차} --- 원 원형은 점유 $0/1$ 만으로
미시상태를 지정하나(내부 자유도 없음), 본문은 점유 자리에 진동 등 내부 자유도 $q(T)$(식~\eqref{eq:partfn})를
실어 유효 자리에너지 $\tilde\varepsilon(T){=}\varepsilon_0{-}\kB T\ln q(T)$ 로 온도 이동을 허용하고(Part T 의
엔트로피$\cdot$발열 연결), 상호작용도 규칙용액 $\Omega_j$ 로 파라미터화해 spinodal 문턱 $\Omega_j{>}2RT$ 를
명시 분기로 닫는다.
\end{srcbox}

\begin{verifybox}
극한 검산 --- (i) $\beta\to\infty$($T\to0$): $\varepsilon_0\gtrless\mu$ 에 따라 $\langle n\rangle^{0}\to0/1$
의 계단 --- $E_i-\mu N_i$ 가 낮은 상태만 남는다는 \S\ref{sec:sm-ensemble} 의 검산과 합치한다.
(ii) $\beta\to0$: $\langle n\rangle^{0}\to\tfrac12$ --- 두 미시상태 등확률. (iii) $\mu\to\pm\infty$:
$\langle n\rangle^{0}\to1/0$ --- 저장조가 이온을 한없이 밀어 넣는/빼내는 극한. 부호 읽기 --- 분모의 지수가
$+\beta(\varepsilon_0-\mu)$ 이므로 점유가 유리할수록($\varepsilon_0<\mu$) 지수가 음으로 죽어
$\langle n\rangle^{0}$ 이 $1$ 에 붙는 방향이다.
\end{verifybox}

\begin{bgbox}
\textbf{페르미온·보손과 양자 통계 ---} 양자역학에서 같은 종류의 입자 둘은
\emph{동일입자}(identical particles)다 --- 전자 둘을 맞바꿔도 어떤 측정값도 달라지지 않는다. 이
구별불가능성은 다입자 파동함수에 교환 대칭 제약을 건다: 두 입자의 자리를 맞바꿀 때
\begin{equation}
\psi(2,1)\;=\;\pm\,\psi(1,2)\qquad(+:\ \text{대칭 = \emph{보손}},\quad -:\ \text{반대칭 = \emph{페르미온}}),
\label{eq:sm-exch}
\end{equation}
의 두 부류만 자연에 존재한다(정수 스핀 $\leftrightarrow$ 대칭, 반정수 스핀 $\leftrightarrow$ 반대칭의
대응이 \emph{스핀--통계 정리}이고 --- 광자$\cdot$포논이 보손(boson), 전자가 페르미온(fermion)의 예다).
반대칭 부류에서 같은 상태에 둘을 넣으면 식~\eqref{eq:sm-exch} 이 $\psi(1,1)=-\psi(1,1)$, 곧 $\psi\equiv0$
을 강제하므로 준위당 점유가 $0$ 또는 $1$ 로 제한된다(\emph{Pauli 배타 원리}, Pauli exclusion principle);
대칭 부류에는 이 제약이 없어 한 상태를 몇 개라도 함께 점유한다. 두 부류의 평형 점유 분포는 본문과 같은
대정준 도구(식~\eqref{eq:sm-gc})로 준위 하나에서 곧장 닫힌다 --- 준위 $E$ 의 점유수 합이 페르미온은
$n\in\{0,1\}$ 두 항, 보손은 $n=0,1,2,\dots$ 의 기하급수라($E>\mu$ 에서 수렴)
\begin{equation}
\begin{aligned}
\text{페르미온: }&\ \Xi_\mathrm F=\!\!\sum_{n=0,1}\!e^{-n\beta(E-\mu)}=1+e^{-\beta(E-\mu)}
&&\Rightarrow\ \
f(E)=\frac{1}{e^{\beta(E-\mu)}+1},\\
\text{보손: }&\ \Xi_\mathrm B=\sum_{n=0}^{\infty}e^{-n\beta(E-\mu)}=\frac{1}{1-e^{-\beta(E-\mu)}}
&&\Rightarrow\ \
\bar n(E)=\frac{1}{e^{\beta(E-\mu)}-1}
\end{aligned}
\label{eq:sm-fdbe}
\end{equation}
이고, 두 화살표는 같은 한 줄 $\langle n\rangle=k_BT\,\partial\ln\Xi/\partial\mu$(식~\eqref{eq:sm-gc} 셋째
식)로 나온다. 앞이 Fermi--Dirac 분포, 뒤가 Bose--Einstein 분포다 \cite{mcquarrie1976}. 격자 삽입 자리의
점유 $0/1$ 은 페르미온과 이 셈법 구조만 공유한다 --- $\Xi_\mathrm F$ 와 원형~\eqref{eq:sm-baresum} 의 두-항
합이 같은 대수라 같은 logistic(Fermi--Dirac 형) 분포가 나오지만, 배타의 근거가 파동함수의 반대칭성이 아니라
정전$\cdot$입체 반발이라는 \emph{기하}이고 입자도 전자가 아니라 Li 이온이다. 실제 양자 통계가 일하는 곳은
Part T(\S\ref{sec:vibel})다 --- 격자 진동의 양자(포논)는 보손이라 진동 엔트로피가 식~\eqref{eq:sm-fdbe} 의 $\bar n$
(Bose--Einstein 들뜸수)으로 적히고, 전도 전자는 페르미온이라 electronic 엔트로피가 $f(E)$(Fermi--Dirac
분포)로 적힌다 \cite{mcquarrie1976,ashcroftmermin1976}.
\end{bgbox}

\textbf{(a$'$) 출발 --- 미시상태의 온전한 기술.} 원형의 점입자 이상화는 실제 자리에서 성립하지 않는다 ---
점유수 $n$ 만으로는 미시상태가 다 지정되지 않기 때문이다. $n=1$ 인 자리에서 이온은 한 점에 얼어붙어
있는 것이 아니라 자리 우물 안에서 진동하는 연속 자유도(위치$\cdot$운동량)를 가진다. 따라서 미시상태의
목록을 ``빈 자리($n=0$, 에너지 $0$ 기준)'' 하나와, ``점유 자리($n=1$, 우물 바닥 에너지
$\varepsilon_0$)에 내부 배치 $\Gamma$ 가 붙은 연속족''으로 넓힌다 --- 가정 1$\cdot$2 는 그대로 둔다.

\textbf{(b$'$) 연산 --- 점유수 합 $\times$ 내부 자유도 적분.} 식~\eqref{eq:sm-gc} 의 대정준 합을 이 목록
그대로 쓰면, 점유수에 대한 합(가정 1 로 두 항)과 각 점유 상태의 내부 배치에 대한 적분으로 갈라진다:
\begin{equation}
\Xi_1\;=\;\sum_{n=0,1}e^{\beta\mu n}\,z_n
\;=\;1+q(T)\,e^{-\beta(\varepsilon_0-\mu)},
\qquad
q(T)\equiv\int e^{-\beta H_\mathrm{int}(\Gamma)}\,\frac{\dd\Gamma}{h^3}
\label{eq:partfn}
\end{equation}
--- $z_n$ 은 점유수 $n$ 인 상태들의 canonical 분배함수로, $z_0=1$(빈 자리는 상태 하나),
$z_1=e^{-\beta\varepsilon_0}q(T)$(우물 바닥 인자에 내부 자유도의 적분 $q(T)$ 가 곱해진다)이다.
$q(T)$ 는 점유된 이온의 진동 등 국소 자유도의 단일 자리 분배함수다\footnote{이 $q(T)$(단일 자리 내부 분배함수)는 N6--N9 의 정규화 용량 좌표 $q{=}Q/Q_\cell$ 와 다른 양이다 --- 전자는 Part 0 의 국소 열역학량, 후자는 곡선 좌표.}. 일반적인 정의는 상태
합 $q(T)=\mathrm{Tr}\,e^{-\beta H_\mathrm{int}}$ 이고 식~\eqref{eq:partfn} 의 위상공간 적분은 그 고전
극한이다 --- 자리 우물을 3차원 조화 퍼텐셜(진동수 $\omega_0$)로 근사하면 고전 극한에서
$q=(k_BT/\hbar\omega_0)^3$ 이고, 등방 가정(세 방향이 같은 $\omega_0$)을 풀어 데카르트 축마다 다른
진동수를 허용하는 일반형으로 넓히면 양자 상태 합으로는
$q=\prod_{i=1}^{3}[2\sinh(\beta\hbar\omega_i/2)]^{-1}$(축 지표 $i=1,2,3$ 마다 진동수 $\omega_i$)로 닫힌다.
내부 분배함수 $q(T)$ 를 포함한 이 확장 역시 국소화 흡착 통계의 표준 전개이고
\cite{hill1960,fowler1939,mcquarrie1976}, $q\equiv1$(내부 자유도 없음)로 두면
식~\eqref{eq:partfn} 는 원형~\eqref{eq:sm-baresum} 으로 되돌아간다. 위 첨자 없는 $\Xi_1$ 은 내부
자유도까지 포함한 완전한 단일 자리 대정준 분배함수이고, Part T 의 단일 자리 분배함수(\S\ref{ssec:site},
식~\eqref{eq:Z1})도 같은 기호 $\Xi_1$ 로 통일한다.

\textbf{(c$'$) 중간식 --- 평균 점유와 계수 $q(T)$ 의 흡수.} 점유 상태들의 확률 합이 곧 평균 점유다:
\begin{equation}
\langle n\rangle=0\cdot p_0+1\cdot p_1
=\frac{q(T)\,e^{-\beta(\varepsilon_0-\mu)}}{1+q(T)\,e^{-\beta(\varepsilon_0-\mu)}},
\label{eq:sm-occmid}
\end{equation}
그리고 식~\eqref{eq:sm-gc} 의 셋째 식으로 교차검증하면 $k_BT\,\partial\ln\Xi_1/\partial\mu$ 가 같은
결과를 준다. 여기서 적분이 남긴 계수 $q(T)$ 는 \emph{분자와 분모의 점유-상태 항 양쪽에 같은 인수로}
붙어 있다. 따라서 결과는 $q$ 와 $\varepsilon_0$ 을 따로 아는 데 의존하지 않고 한 조합에만 의존하며,
그 조합은 지수 안으로 흡수된다:
\begin{equation}
q(T)\,e^{-\beta\varepsilon_0}\;=\;e^{-\beta\tilde\varepsilon(T)},
\qquad
\tilde\varepsilon(T)\;\equiv\;\varepsilon_0-k_BT\ln q(T)
\label{eq:sm-epstilde}
\end{equation}
--- $\tilde\varepsilon$ 는 점유 자리의 \emph{자유에너지}(우물 바닥 에너지 $+$ 내부 자유도의 자유에너지
$-k_BT\ln q$)다. 자리 점유의 자유에너지 비용을 $\Delta\mu\equiv\tilde\varepsilon-\mu_\mathrm{Li}$ 로
줄여 쓰면 식~\eqref{eq:sm-occmid} 는 $e^{-\beta\Delta\mu}/(1+e^{-\beta\Delta\mu})$ 가 되어, 내부
자유도가 있어도 함수형은 원형 박스~\eqref{eq:sm-bare} 의 두-상태 logistic 그대로 보존된다.

\textbf{(d$'$) 박스 --- 점유율 $\theta$.} 분모·분자를 $e^{-\beta\Delta\mu}$ 로 나누면
\begin{equation}
\boxed{\;\theta\;\equiv\;\langle n\rangle\;=\;\frac{1}{1+e^{+\beta\,\Delta\mu}}\;,
\qquad \Delta\mu\equiv\tilde\varepsilon(T)-\mu_\mathrm{Li}\;.}
\label{eq:fermifn}
\end{equation}
\begin{verifybox}
극한 검산 --- (i) $\beta\to\infty$($T\to0$)에서 $\Delta\mu\gtrless0$ 에 따라 $\theta\to0/1$ 의 계단
($T\to0$ 에서 $\tilde\varepsilon$ 는 점유 상태의 바닥 에너지로 수렴한다 --- 고전 $q$ 면 $k_BT\ln q\to0$
이라 $\varepsilon_0$, 양자 조화 $q$ 면 영점 에너지를 더한 $\varepsilon_0+\tfrac32\hbar\omega_0$ --- 어느
쪽이든 상수라 계단 구조는 동일). (ii) $\beta\to0$($T\to\infty$)에서는 $\beta\Delta\mu\to-\ln q$ 이므로
내부 자유도가 없는 기준($q\equiv1$)에서 $\theta\to\tfrac12$(두 상태 등확률)이고, $q(T)$ 가 온도와 함께
증가하는 실제 조화 우물에서는 점유 상태의 위상공간 증가가 이겨 $\theta\to1$ 로 향한다.
(iii) $\mu_\mathrm{Li}\to\pm\infty$ 에서 $\theta\to1/0$. (iv) 확장을 끄는 극한 $q\equiv1$ 에서
$\tilde\varepsilon\to\varepsilon_0$·$\Delta\mu\to\varepsilon_0-\mu_\mathrm{Li}$ 가 되어($\mu_\mathrm{Li}$ 는 원형의 $\mu$ 와
같은 저장조 화학퍼텐셜의 자리당 표기) 식~\eqref{eq:fermifn} 는 원형 박스~\eqref{eq:sm-bare} 의 $\langle n\rangle^{0}$ 을 그대로 회수한다.
★함수형 가드 --- 식~\eqref{eq:fermifn} 은 Fermi--Dirac 함수와 \emph{동형}이지만 물리량은 다르다:
공유하는 것은 ``한 자리에 입자 0 또는 1''이라는 배타 점유의 대수 구조뿐이고(가정 1), 여기의 입자는
전자가 아니라 Li 이온이며 배타성의 근거도 페르미온의 Pauli 배타 같은 양자 통계가 아니라 자리 하나에
이온 하나라는 기하다.
\end{verifybox}

한 걸음 더 --- $n\in\{0,1\}$ 라 $n^2=n$ 이므로 점유의 fluctuation(분산)이 닫힌 꼴로 나온다:
\begin{equation}
\mathrm{var}(n)=\langle n^2\rangle-\langle n\rangle^2=\theta(1-\theta),
\qquad
\frac{\partial\theta}{\partial(\beta\mu)}=\theta(1-\theta)
\label{eq:sm-flucres}
\end{equation}
(둘째 식은 식~\eqref{eq:fermifn} 의 $\beta\mu$ 미분). \emph{점유의 $\mu$-감도 $=$ 점유 fluctuation} --- fluctuation--response 관계의 최소 사례이고, 본론 평형 peak 의 종 $\xi_\eq(1-\xi_\eq)$(\S\ref{sec:eqpeak})의
통계적 정체가 바로 이 분산이다.

유효 자리 자유에너지 $\tilde\varepsilon(T)$ 가 하류에서 가는 곳은 두 갈래다. 첫째, 온도를 고정하고 보면
$\tilde\varepsilon$ 는 상수 하나이고, \S\ref{sec:sm-electro} 의 전기화학 결선에서 전이 중심 $U_j$ 로
흡수된다 --- 그 소절과 본론이 쓰는 몰당 기준 에너지 $\mu^0$ 는 정확히 $\tilde\varepsilon$ 의 몰 환산
$N_A\tilde\varepsilon$ 이다. 둘째, 온도를 미분하면 내부 자유도의 \emph{자리당} 자유에너지
$f_\mathrm{int}=-k_BT\ln q$(소문자 $f$ --- 계 전체 Helmholtz $F{=}-k_BT\ln Z$ 가 아니라 자리당 양) 가 자리당 엔트로피
\begin{equation}
s_\mathrm{int}\;=\;-\frac{\partial f_\mathrm{int}}{\partial T}
\;=\;k_B\,\frac{\partial\,(T\ln q)}{\partial T}
\label{eq:sm-sint}
\end{equation}
를 주는데, 이것이 Part T 의 vibrational 엔트로피 사슬(\S\ref{sec:vibel})의 단일 자리 원형이다 --- 여기 자리 하나의
세 축 진동수 $\omega_i$($i=1,2,3$)를 격자 전체의 모든 진동 모드를 헤아리는 지표 $k$ 로 넓혀, 조화 모드의
$q_k=[1-e^{-\beta\hbar\omega_k}]^{-1}$ 를 식~\eqref{eq:sm-sint} 에 넣으면 Part T 의 모드당 엔트로피~\eqref{eq:Svib_mode}
$s_k=k_B[(1+\bar n_k)\ln(1+\bar n_k)-\bar n_k\ln\bar n_k]$(Bose--Einstein 들뜸수 $\bar n_k$)가 그대로
나온다. 기준점에 관한 사소한 주의로, 위 (b$'$) 의 $[2\sinh(\beta\hbar\omega/2)]^{-1}$ 은 우물 바닥 기준
(영점 에너지를 $q$ 에 포함)이고 여기의 $[1-e^{-\beta\hbar\omega}]^{-1}$ 은 영점 에너지를
$\varepsilon_0$ 쪽에 포함시킨 기준이다 --- 두 규약의 $q$ 는 인자 $e^{-\beta\hbar\omega/2}$(영점 에너지
$\hbar\omega/2$ 의 Boltzmann 인자)만큼 다른데, 이 인자가 $T\ln q$ 에 남기는 기여는 온도 무관 상수
$-\hbar\omega/(2k_B)$ 이므로 온도 미분, 곧 엔트로피는 동일하다. 따라서 전이 중심의 온도 기울기
$\dd U_j/\dd T=\Delta S_{\rxn,j}/F$(\S\ref{sec:center})에서 진동
엔트로피가 차지하는 몫의 미시적 기원이 바로 $\tilde\varepsilon(T)$ 의 온도 의존, 곧 $q(T)$ 다.

자리 하나가 준 답 --- 점유율 $\theta$, fluctuation $\theta(1-\theta)$, 유효 자리 자유에너지
$\tilde\varepsilon(T)$ --- 을 손에 쥐었으므로, 다음 소절은 같은 답을 $M$ 개의 독립 자리로 복제해 거시
점유율과 섞임 엔트로피로 넓힌다.

\subsection{$M$ 개의 독립 자리 --- lattice gas 와 섞임 엔트로피}\label{sec:sm-lattice}
\textbf{(a) 출발 --- 동등·독립 자리의 집합.} 결정 안에 동등한 삽입 자리가 $M$ 개 있고 서로 상호작용하지
않는다고 가정한다 --- 이 모형이 \emph{격자기체}(lattice gas)의 이상 극한이다. 미시상태는 점유 배열
$(n_1,\dots,n_M)$($n_k\in\{0,1\}$, 가정 1)에 각 점유 자리의 내부 배치가 붙은 것이고, $N=\sum_kn_k$ 다.

\textbf{(b) 연산 --- 분배함수의 인수분해.} 자리들이 독립(가정 2)이라 대정준 합이 자리별 곱으로 갈라지고,
각 인수는 내부 자유도 적분까지 마친 단일 자리 결과~\eqref{eq:partfn} 그대로다
($\Delta\mu=\tilde\varepsilon-\mu$ 로 흡수된 표기):
\begin{equation}
\Xi_M=\sum_{n_1=0,1}\!\cdots\!\sum_{n_M=0,1}\prod_{k=1}^{M}e^{-\beta\,\Delta\mu\,n_k}
=\prod_{k=1}^{M}\Big[\sum_{n_k=0,1}e^{-\beta\,\Delta\mu\,n_k}\Big]=\Xi_1^{\,M},
\label{eq:sm-factor}
\end{equation}
따라서 식~\eqref{eq:sm-gc} 의 셋째 식으로 $\langle N\rangle=k_BT\,\partial(M\ln\Xi_1)/\partial\mu=M\theta$ ---
거시 점유율이 단일 자리 점유 확률 \eqref{eq:fermifn} 와 같다.

\textbf{(c) 중간식 --- 같은 결과의 셈법 경로(교차검증).} 이번엔 입자수 $N$ 을 고정하고 배치를 센다.
$M$ 자리에서 $N$ 자리를 고르는 경우의 수 $W=M!/[N!(M-N)!]$ 에 공준 2(식~\eqref{eq:sm-S})와 Stirling 근사
$\ln m!\simeq m\ln m-m$ 을 적용하면($\theta=N/M$)
\begin{equation}
S_\mathrm{mix}=k_B\ln W\simeq-k_BM\big[\theta\ln\theta+(1-\theta)\ln(1-\theta)\big],
\label{eq:sm-Smix}
\end{equation}
자유에너지 $F(N)=N\tilde\varepsilon-TS_\mathrm{mix}$(점유 자리당 자유에너지는 내부 자유도 몫까지 포함한
$\tilde\varepsilon$, 식~\eqref{eq:sm-epstilde})를 $N$ 으로 미분하면
\begin{equation}
\mu=\frac{\partial F}{\partial N}
=\tilde\varepsilon+k_BT\ln\frac{\theta}{1-\theta}.
\label{eq:sm-mucount}
\end{equation}
한편 식~\eqref{eq:fermifn} 를 $\mu$ 에 대해 뒤집어도($e^{\beta\Delta\mu}=(1-\theta)/\theta$) 같은
식~\eqref{eq:sm-mucount} 가 나온다 --- 대정준 경로(자리 하나의 확률)와 셈법 경로(배치 수의 로그)가 같은
$\mu(\theta)$ 에 닿는 것이 앙상블 동등성이며, 본론 \S\ref{sec:dist} 가 말하는 ``kinetic·thermo 두 경로가 한
분포의 두 측면''의 평형 쪽 절반이 바로 이것이다.

\textbf{(d) 박스 --- 몰 단위 이상 lattice gas 화학퍼텐셜.} 자리당 $k_B$·$\tilde\varepsilon$ 을 몰당
$R=N_Ak_B$·$\mu^0\equiv N_A\tilde\varepsilon$ 으로 환산하면
\begin{equation}
\boxed{\;\mu(\theta)=\mu^0+RT\ln\frac{\theta}{1-\theta}\qquad(\text{이상 lattice gas})\;.}
\label{eq:sm-muideal}
\end{equation}
\begin{verifybox}
부호·극한 검산 --- $\theta\to0$ 에서 $\mu\to-\infty$(거의 빈 격자에 입자를 넣는 일은 섞임 엔트로피 이득 때문에
얼마든지 값싸다), $\theta\to1$ 에서 $+\infty$(마지막 빈자리는 무한히 비싸다), $\theta=\tfrac12$ 에서
$\mu=\mu^0$·자리당 $S_\mathrm{mix}$ 최대 $k_B\ln2$. 식~\eqref{eq:sm-muideal} 는 다음 소절 $\mu(\theta)$ 의
$\Omega=0$ 특수형이고, 상호작용 몫 $\Omega(1-2\theta)$ 를 다음 소절이 더한다.
\end{verifybox}

% 자산(v1.0.20 신규): [V20-001 원형 2-상태 선행 유도 eq:sm-bare] [V20-002 페르미온/보손 배경 bgbox]
% 자산: [A-015] [A-016] [A-017] [A-018] [A-019] [A-020] [A-021] [A-022] [A-023] [A-024] [A-025] [A-026] [A-027] [A-028] [A-029] [A-030] [A-031] [A-032] [A-033] [A-034] [A-035] [A-036] [A-037] [A-038] [A-039] [A-040] [A-041] [A-042] [A-043] [A-044] [A-045] [A-046]
